\section{Arquitectura de la aplicación}

\subsection{Tecnologías}

\begin{longtable}{p{0.22\textwidth} p{0.22\textwidth} p{0.16\textwidth} p{0.30\textwidth}}
\toprule
\textbf{Capa} & \textbf{Tecnología} & \textbf{Versión} & \textbf{Propósito} \\
\midrule
Interfaz & React & 19 & Construcción basada en componentes \\
Lenguaje & TypeScript & 5.9 & Seguridad de tipos \\
Construcción & Vite & 7 & Desarrollo y compilación de producción \\
Enrutamiento & React Router & 7 & Navegación entre pantallas \\
Estilos & CSS nativo & Actual & Diseño responsive y temas visuales \\
\bottomrule
\end{longtable}

\subsection{Organización}

La aplicación separa las rutas en un archivo principal y concentra las pantallas en el módulo de páginas. El sistema visual global se mantiene en una hoja CSS compartida. La documentación técnica se conserva fuera de \texttt{src} para no mezclar código ejecutable con artefactos de análisis.

\begin{itemize}[leftmargin=*]
  \item \textbf{src/App.tsx}: definición central de rutas.
  \item \textbf{src/pages/Pages.tsx}: layout y pantallas del frontend.
  \item \textbf{src/styles.css}: estilos responsive, interacción y temas.
  \item \textbf{docs}: índice y bitácora de desarrollo.
\end{itemize}

\subsection{Servicios futuros}

La integración futura deberá incorporar una API para persistir perfiles, verificaciones, cuentas bancarias, ofertas, cobros, entregas y notificaciones. También deberá definirse un almacenamiento seguro para los documentos y un proveedor de mapas para las ubicaciones.
